% Torsion of a circular shaft — analytical derivation
% AIM Interactive Mechanics series. Compile with pdflatex.
\documentclass[10pt,a4paper]{article}
\usepackage[margin=2cm]{geometry}
\usepackage{amsmath}
\usepackage[T1]{fontenc}
\usepackage{lmodern}
\usepackage{microtype}
\usepackage{titlesec}
\titlespacing*{\section}{0pt}{9pt}{2pt}
\usepackage{xcolor}
\definecolor{iron}{HTML}{935636}
\definecolor{celdark}{HTML}{4B6C61}
\definecolor{celline}{HTML}{CDE0D8}
\usepackage[colorlinks=true,urlcolor=iron,linkcolor=black]{hyperref}
\setlength{\parindent}{0pt}
\setlength{\parskip}{4pt}
\AtBeginDocument{\setlength{\abovedisplayskip}{6pt}\setlength{\belowdisplayskip}{6pt}\setlength{\abovedisplayshortskip}{3pt}\setlength{\belowdisplayshortskip}{3pt}}
\pagestyle{empty}

\begin{document}

{\footnotesize\textcolor{celdark}{\textsc{aim interactive mechanics \,\textperiodcentered\, analytical derivation}}}

{\LARGE Torsion of a Circular Shaft}

\textcolor{celdark}{\small Companion to the interactive simulator at
\url{https://aimlab.uk/sims/torsion.html}}

\medskip
{\color{celline}\hrule height 1pt}

\section*{Setup and assumptions}
A solid circular shaft of diameter $d$ and length $L$ is fixed at $x = 0$ and
twisted by a torque $T$ at $x = L$. The material is linear elastic with shear
modulus $G$; rotations are small. Let $\varphi(x)$ be the rotation of the
cross-section at $x$.

\section*{Kinematics}
By circular symmetry, every cross-section rotates rigidly in its own plane ---
plane sections remain plane and undistorted. This is true only for circular sections;
any other shape warps. Consider two neighbouring sections a distance
$\mathrm{d}x$ apart, rotated relative to each other by $\mathrm{d}\varphi$.
A material line of length $\mathrm{d}x$ on the cylinder of radius $r$ tilts by
the angle
\begin{equation}
\gamma(r) \;=\; \frac{r \,\mathrm{d}\varphi}{\mathrm{d}x} \;=\; r\,\varphi' ,
\end{equation}
which is the shear strain: zero on the axis, largest at the surface.

\section*{Constitutive law and stress resultant}
Hooke's law in shear gives $\tau(r) = G\gamma(r) = G\,r\,\varphi'$. The torque
carried by the section is the moment of this stress distribution:
\begin{equation}
T(x) \;=\; \int_A \tau\, r \,\mathrm{d}A
\;=\; G\varphi' \int_A r^2 \,\mathrm{d}A
\;=\; G J \varphi',
\qquad
J \;=\; \int_0^{d/2} r^2 \,(2\pi r)\,\mathrm{d}r \;=\; \frac{\pi d^4}{32}.
\end{equation}
$J$ is the polar second moment of area --- the torsional analogue of $I$.

\section*{Governing equation}
With no distributed torque along the shaft, equilibrium of a slice gives
$\mathrm{d}T/\mathrm{d}x = 0$, so
\begin{equation}
G J \,\frac{\mathrm{d}^2 \varphi}{\mathrm{d}x^2} \;=\; 0,
\qquad
\varphi(0) = 0, \qquad G J\,\varphi'(L) = T .
\end{equation}
Exactly the same second-order structure as the axial rod, with
$(u, EA, F) \rightarrow (\varphi, GJ, T)$.

\section*{Solution}
\begin{equation}
\boxed{\; \varphi(x) \;=\; \frac{T x}{G J} \;}
\qquad\Rightarrow\qquad
\varphi_L \;=\; \frac{T L}{G J},
\qquad
\tau_{\max} \;=\; \frac{T\,(d/2)}{J} \;=\; \frac{16\,T}{\pi d^3},
\qquad
\gamma_{\max} \;=\; \frac{(d/2)\,\varphi_L}{L}.
\end{equation}
Note that the twist depends on $G$ but the stress does not: change the material
in the simulator and watch $\varphi$ collapse while $\tau_{\max}$ stays put.
Since $J \propto d^4$, doubling the diameter makes the shaft sixteen times
stiffer in twist.

\section*{Worked example (simulator defaults)}
$G = 10$~MPa, $d = 20$~mm, $L = 300$~mm, $T = 0.5$~N\,m:
\begin{equation}
J = \frac{\pi \times 20^4}{32} = 15\,708~\text{mm}^4, \qquad
\varphi_L = \frac{0.5 \times 0.3}{10^7 \times 1.571\times10^{-8}}
= 0.955~\text{rad} = 54.7^\circ,
\end{equation}
with $\tau_{\max} = 0.318$~MPa and $\gamma_{\max} = 3.18\%$. Set the simulator
sliders to these values and compare the dial. Beyond roughly $10\%$ surface shear
strain the small-strain assumption becomes invalid.

\vfill
{\color{celline}\hrule height 0.6pt}
\smallskip
{\footnotesize\textcolor{celdark}{Part of the multimodal mechanics series
(analytical \textperiodcentered\ simulated \textperiodcentered\ experimental)
\,\textperiodcentered\, AIM Group, Imperial College London.}}

\end{document}
